\section{Results}
\label{sec:results}

\subsection{Estimator verification}
\label{sec:estimator-results}

The local equivalence tier first asks whether the scalable randomized adjoint estimates the same
transported-defect profile as exact discrete-Heun JVP propagation. With $R=4$, normalized profile
error remains below the declared 10\% threshold on five modes, fan modes, spiral, GM16 seed 0, and
hard Burgers. The resulting integration error remains within 5\% of the JVP schedule, and varying
the adjoint seed among 7, 2718, and 31415 changes endpoint integration error by less than 5\%.
These checks support the estimator, not the optimality of the resulting mesh.

The staged-shapes image pilot is a useful stress test of that conclusion. Its $R=4$ adjoint
profile differs from JVP by only $0.020\%$ in normalized L2, and the resulting integration error
differs by at most $1.47\%$ across the four NFE budgets. Nevertheless, the maximum knot
displacement is 0.034, above the declared 0.02 gate, because the profile is sharply concentrated
near zero and small profile perturbations shift cumulative-density quantiles. Probe-seed variation
passes at NFE 5, 10, and 50 but reaches $5.06\%$ at NFE 20, narrowly missing the 5\% gate. We
retain both failures: randomized adjoints estimate the transported profile accurately here, but
mesh inversion can amplify small errors and should be stabilized in future implementations.

Both CIFAR-10 pilots pass all declared estimator gates. For conditional and unconditional models,
respectively, the $R=4$ profile errors are 1.56\% and 1.17\%, and maximum knot displacements are
0.0066 and 0.0082. Their integration errors differ from exact JVP schedules by at most 0.30\%,
while variation over the three probe seeds remains below 0.70\%. The finite-difference profiles
also pass both perturbation checks. Thus the negative CIFAR endpoint result below is not explained
by a failed randomized estimator or unstable knot construction.

Calibration cost is reported outside sampling NFE. On the A100 used for image evaluation, the
local CIFAR profile takes approximately 15 seconds; the four-probe adjoint profile takes 246--255
seconds and peaks below 0.85 GB of allocated GPU memory. Amortized over the reported 5000 samples,
the latter is about 0.05 seconds per sample. This cost is paid once per checkpoint and conditioning
protocol, so the method is most appropriate when the resulting population mesh is reused. Exact
JVP calibration takes approximately 1772--1779 seconds on the same pilots, about seven times the
$R=4$ adjoint cost.

Centered finite differences agree with JVP at the declared perturbations on fan, spiral, GM16, and
Burgers. Five modes is a retained failure: $\varepsilon\in\{10^{-3},3\times10^{-3}\}$ lies outside
the local linear regime and produces normalized profile errors of 14.1\% and 27.6\%. A separate
convergence audit reaches 0.23\% at $\varepsilon=10^{-4}$ before float32 cancellation becomes
visible. Staged-shapes exposes the complementary issue: $\varepsilon=10^{-3}$ gives 0.92\%
profile error, yet a 0.123 maximum knot displacement, while
$3\times10^{-3}$ fails both profile and knot gates. We do not replace either declared result with
a favorable post-hoc perturbation.

\begin{figure}[t]
  \centering
  \safeincludegraphics[width=\linewidth]{../results/sctw_adjoint_suite/aggregate/profile_comparison_nfe20.png}
  \caption{Population SCTW and Adjoint-SCTW profiles. Each curve is normalized to unit integral;
  colored rugs show the corresponding $p=0.5$, NFE-20 knots. Terminal sensitivity can amplify
  early defects, as in Burgers and staged-shapes, or damp a sharp local endpoint profile, as in
  conditional CIFAR-10.}
  \label{fig:profile-comparison}
\end{figure}

\subsection{When integration accuracy predicts endpoint quality}
\label{sec:alignment-results}

We measure learned-ODE integration error against the declared fine Heun reference and separately
measure endpoint quality. Their relationship is benchmark dependent. Across the 14 common
schedules, staged-shapes feature W1 and integration error are almost perfectly rank-aligned:
Spearman correlation ranges from 0.969 to 0.996 across the four NFE budgets. Conditional CIFAR-10
is weakly or negatively aligned from NFE 10 onward, reaching $-0.604$ at NFE 50; unconditional
CIFAR-10 is near zero at NFE 20 and $-0.110$ at NFE 50. These are descriptive within-checkpoint
correlations, not significance tests over independent solvers. They nevertheless identify an
important boundary: better numerical integration helps endpoint quality when the fine learned-ODE
endpoint is itself aligned with the data metric, but cannot correct a biased learned field.

\subsection{Staged-shapes}
\label{sec:staged-results}

\begin{table}[t]
  \centering
  \caption{Staged-shapes easy benchmark; feature W1 is lower better. Best fixed is the
  post-hoc minimum over the six early/late/symmetric power meshes and two fixed VSFM meshes at each
  NFE, so it is a deliberately favorable comparator rather than one deployable preselected mesh.}
  \label{tab:staged}
  \resizebox{\linewidth}{!}{%
    \input{../results/sctw_adjoint_suite/aggregate/quality_staged_shapes.tex}%
  }
\end{table}

\begin{figure}[t]
  \centering
  \begin{subfigure}{0.31\linewidth}
    \centering
    \safeincludegraphics[width=\linewidth]{../results/sctw_adjoint_suite/remote/staged_shapes/samples/uniform_nfe_5.png}
    \caption{Uniform Euler.}
  \end{subfigure}
  \begin{subfigure}{0.31\linewidth}
    \centering
    \safeincludegraphics[width=\linewidth]{../results/sctw_adjoint_suite/remote/staged_shapes/samples/power_early_rho3_nfe_5.png}
    \caption{Best fixed mesh.}
  \end{subfigure}
  \begin{subfigure}{0.31\linewidth}
    \centering
    \safeincludegraphics[width=\linewidth]{../results/sctw_adjoint_suite/remote/staged_shapes/samples/adjoint_sctw_p0p5_nfe_5.png}
    \caption{Adjoint-SCTW.}
  \end{subfigure}
  \caption{Fixed-prior NFE-5 staged-shapes samples. Adjoint-SCTW removes more incomplete and
  ambiguous shapes than uniform Euler and slightly improves over the strongest fixed mesh.}
  \label{fig:staged-samples}
\end{figure}

Staged-shapes is the clearest positive endpoint result. At NFE 5, feature W1 falls from 1.524 for
uniform Euler to 1.017 for SCTW and 0.870 for Adjoint-SCTW; mode hit rate rises from 81.3\% to
94.4\% and 97.4\%, respectively. Adjoint-SCTW also beats the post-hoc best fixed mesh at every NFE.
The corresponding integration error falls from 7.750 to 5.039 and 3.589, so the endpoint and
learned-ODE metrics tell the same story.

\subsection{Low-dimensional mixtures and GM16}
\label{sec:lowdim-results}

\begin{figure}[t]
  \centering
  \begin{subfigure}{0.31\linewidth}
    \centering
    \safeincludegraphics[width=\linewidth]{../results/lowdim_sctw_complementary/formative_panels/five_modes_diagnostic_nfe20.png}
    \caption{Five modes.}
  \end{subfigure}
  \begin{subfigure}{0.31\linewidth}
    \centering
    \safeincludegraphics[width=\linewidth]{../results/lowdim_sctw_complementary/formative_panels/fan_modes_diagnostic_nfe20.png}
    \caption{Fan modes.}
  \end{subfigure}
  \begin{subfigure}{0.31\linewidth}
    \centering
    \safeincludegraphics[width=\linewidth]{../results/lowdim_sctw_complementary/formative_panels/spiral_diagnostic_nfe20.png}
    \caption{Spiral.}
  \end{subfigure}
  \caption{Visual low-dimensional transports. These panels expose the target geometry and knot
  allocation; the quantitative comparison uses fixed priors and the tightened matching metrics.}
  \label{fig:lowdim-panels}
\end{figure}

\begin{table}[t]
  \centering
  \caption{GM16 endpoint W1, averaged over three independently trained checkpoints. Lower is
  better. The table reports both declared tempering powers and the post-hoc best fixed comparator.}
  \label{tab:gm16}
  \resizebox{\linewidth}{!}{%
    \input{../results/sctw_adjoint_suite/aggregate/quality_gm16.tex}%
  }
\end{table}

At NFE 5, learned-ODE integration error is 0.460/0.364/0.372 for
uniform/SCTW/Adjoint-SCTW on five modes, 0.638/0.514/0.479 on fan modes, and
0.275/0.260/0.255 on spiral. Thus terminal sensitivity is not uniformly beneficial: SCTW remains
slightly better on five modes, while Adjoint-SCTW improves fan and spiral. GM16 is stronger and
replicated: the three-seed NFE-5 integration error is
$0.527\pm0.012$, $0.263\pm0.009$, and $0.237\pm0.009$, with corresponding endpoint W1
1.257, 1.143, and 1.128.

\subsection{Hard Burgers solution map}
\label{sec:burgers-results}

\begin{figure}[t]
  \centering
  \safeincludegraphics[width=0.92\linewidth]{../results/burgers_solution_map_dedalus_rich/burgers_data_breadth_examples.png}
  \caption{Examples from the rich one-dimensional Burgers solution-map cache. Gray curves are
  initial conditions and blue curves are final PDE reference solutions.}
  \label{fig:burgers-data}
\end{figure}

\begin{figure}[t]
  \centering
  \safeincludegraphics[width=0.92\linewidth]{../results/burgers_solution_map_dedalus_rich/burgers_solution_map_dedalus_rich_h256_seed42/eval_sctw_schedules_ema/profile_diagnostic_nfe20.png}
  \caption{Hard Burgers local-profile diagnostic. The newer transported-profile comparison is
  included in \cref{fig:profile-comparison}; this earlier figure retains the solution trajectories
  and endpoint metrics.}
  \label{fig:burgers-diagnostic}
\end{figure}

\begin{table}[t]
  \centering
  \caption{Hard Burgers integration error to Heun-1000 for the learned ODE; lower is better.}
  \label{tab:burgers-integration}
  \input{../results/sctw_adjoint_suite/aggregate/integration_burgers.tex}
\end{table}

Adjoint-SCTW gives the largest integration gain on Burgers: at NFE 5 the error falls from 1.756
for uniform Euler and 1.673 for SCTW to 1.250, and at NFE 50 it falls to 0.179. Endpoint PDE RMSE,
however, changes only from 0.357 to 0.350 at NFE 5 and is essentially unchanged at higher NFE.
The method integrates the learned ODE more accurately; it does not correct the learned
initial-to-final map.

\subsection{CIFAR-10}
\label{sec:cifar-results}

\begin{table}[t]
  \centering
  \caption{Conditional CIFAR-10 minibatch-OT checkpoint, EMA evaluation with 5000 fixed-prior
  samples. FID is lower better. Results are descriptive for one checkpoint and evaluation seed.}
  \label{tab:cifar-conditional}
  \resizebox{\linewidth}{!}{%
    \input{../results/sctw_adjoint_suite/aggregate/quality_cifar10_conditional.tex}%
  }
\end{table}

\begin{table}[t]
  \centering
  \caption{Unconditional CIFAR-10 Sinkhorn checkpoint under the same 5000-sample protocol.
  FID is lower better; best fixed has the same post-hoc definition as in \cref{tab:staged}.}
  \label{tab:cifar-unconditional}
  \resizebox{\linewidth}{!}{%
    \input{../results/sctw_adjoint_suite/aggregate/quality_cifar10_unconditional.tex}%
  }
\end{table}

SCTW improves FID over uniform Euler at NFE 10, 20, and 50 for both checkpoints. At NFE 5,
tempering matters: $p=0.25$ is slightly better than uniform for conditional CIFAR but neither SCTW
power improves unconditional FID. Adjoint-SCTW consistently lowers learned-ODE integration error
but generally trails SCTW in FID. For conditional CIFAR at NFE 20, integration error is
3.271/3.415/2.662 for uniform/SCTW/Adjoint-SCTW, whereas FID is 23.158/21.391/22.113.
This is a useful negative result for the stronger premise that terminally accurate integration
must improve generative quality.

\subsection{Summary of evidence}
\label{sec:summary-evidence}

Adjoint-SCTW satisfies the stated decision rule: it improves or closely matches SCTW's integration
error on most benchmark families and closes the fixed-mesh gap on the non-Burgers staged-shapes
problem. The result does not justify a universal sampler claim. SCTW is more reliable for CIFAR
FID, and five modes provides a small case where terminal weighting is unnecessary. The combined
evidence supports two distinct claims: local material acceleration is a practical population mesh
signal, and terminal sensitivity can sharpen that signal when numerical integration error is
aligned with the endpoint task.
